\subsection{Datasets}
\label{sec:datasets}

All training data except the epidemic curve are simulated from known
equations, with ground truth generated by SciPy's \cite{virtanen2020scipy}
DOP853 integrator at relative and absolute tolerance $10^{-12}$, far below
any error we report. Table~\ref{tab:datasets} summarises the four systems
and their train/test splits. Every accuracy number is computed only on the
test window, which the training loss never sees.

\paragraph{Lotka-Volterra (main benchmark).} The base paper's predator-prey
system \cite{koenig2024kanode},
\begin{gather}
\dot u_1 = \alpha u_1 - \beta u_1 u_2, \qquad \dot u_2 = \delta u_1 u_2 - \gamma u_2,
\nonumber\\
(\alpha,\beta,\gamma,\delta) = (1.5,\,1.0,\,3.0,\,1.0),
\label{eq:lv}
\end{gather}
from $\mathbf{u}(0)=(1,1)$, sampled every $\Delta t=0.1$. We train on
$t\in[0,3.5]$ (36 points) and score strictly on $t\in(3.5,14.0]$ (105
unseen points). Extrapolation is scored this way because it is the only
meaningful test of whether the network learned the dynamics rather than
memorising the training window; a model that reproduces the 36 training
points can still be badly wrong about the system. Two properties of
Equation~\ref{eq:lv} serve as exact, model-independent checks later on. It
has a single interior equilibrium
$\mathbf{u}^\ast=(\gamma/\delta,\ \alpha/\beta)=(3,\,1.5)$, where the
Jacobian has purely imaginary eigenvalues $\pm i\sqrt{\alpha\gamma}=\pm2.121i$.
And it conserves the first integral
\begin{equation}
H(u_1,u_2)=\delta u_1-\gamma\ln u_1+\beta u_2-\alpha\ln u_2,
\qquad \frac{\dd H}{\dd t}=0,
\label{eq:lv-invariant}
\end{equation}
so every orbit is a closed level curve of $H$ and the phase plane is filled
with a continuous family of neutral cycles. The orbit through $(1,1)$ has
$H_0=2$ and period $T=3.3175$, so the training window covers $1.06$
periods, the test window $3.2$ further periods, and $[0,14]$ about $4.2$.

\paragraph{SIR epidemic model.} The compartmental model splits a fixed
population into susceptible ($S$), infected ($I$) and recovered ($R$)
fractions,
\begin{equation}
\dot S = -\beta S I, \quad \dot I = \beta S I - \gamma I, \quad \dot R = \gamma I,
\label{eq:sir}
\end{equation}
with $(\beta,\gamma)=(0.35,\,0.1)$ and $(S,I,R)(0)=(0.99,0.01,0)$. It has
a built-in algebraic invariant, $S+I+R=1$, because the three terms of the
field cancel when summed. The invariant holds for any parameter values, so a
model that violates it has learned a field with the wrong structure, not
merely an inaccurate one. We simulate $t\in[0,80]$ days at
$\Delta t=0.5$, train on $[0,50]$ (101 points) and test on $(50,80]$.

\paragraph{Damped pendulum.} The angle $\theta$ and angular velocity
$\omega$ evolve under gravity and viscous friction,
\begin{equation}
\dot\theta = \omega, \qquad \dot\omega = -\frac{g}{L}\sin\theta - \mu\,\omega,
\qquad (g/L,\,\mu) = (9.81,\,0.5),
\label{eq:pendulum}
\end{equation}
from $(\theta,\omega)(0)=(2,0)$. Since $\mu>0$, the mechanical energy
$E=\tfrac{1}{2}mL^2\omega^2 + mgL(1-\cos\theta)$ strictly dissipates,
$\dd E/\dd t = -mL^2\mu\,\omega^2 \le 0$. A model can fit the trajectory
well and still violate this monotonic decay, which makes it a genuinely
independent physical check. We simulate $t\in[0,10]$ at $\Delta t=0.05$.
The original recipe trained on $[0,3]$; the fixed recipe
(Section~\ref{sec:crossdomain}) trains on $[0,5]$ (101 points) and tests on
$(5,10]$. The stiffness study (Section~\ref{sec:track-d}) sweeps $\mu$.

\paragraph{Epidemic outbreak curve.} To test whether our SIR fixes rely on
the SIR equations themselves, we also fit a single outbreak wave that no
compartmental ODE generated. The curve is synthetic but built to look like
observed case data: a Gaussian wave of width 14 centred at day 38, skewed by
a $\tanh$ factor so that its infected peak falls on day 45, with $1\%$
Gaussian observation noise,
smoothed by a 7-day moving average. The state is the infected curve together
with a cumulative-recovered curve, both min-max normalized to $[0,1]$. It
spans 120 days; the default split trains on the first 45 days and tests on
the rest.

\begin{table}[t]
\centering
\footnotesize
\setlength{\tabcolsep}{3pt}
\renewcommand{\arraystretch}{1.12}
\caption{Datasets, networks and train/test splits. Times are in the
system's own units (days for the two epidemic systems).}
\label{tab:datasets}
\begin{tabular}{@{}L{1.55cm} L{2.35cm} C{1.95cm} C{1.95cm}@{}}
\toprule
\hdrrow \thh{System} & \thh{Network (params)} & \thh{Train (points)} & \thh{Test (points)}\\
\midrule
Lotka-Volterra & $[2,10,2]$, $G{=}5$ (240) & $[0,3.5]$ (36) & $(3.5,14]$ (105)\\
SIR & $[3,16,3]$, $G{=}8$ (864) & $[0,50]$ (101) & $(50,80]$ (60)\\
Pendulum & $[2,10,2]$, $G{=}8$ (360) & $[0,5]$ (101) & $(5,10]$ (100)\\
Epidemic curve & $[2,10,2]$, $G{=}5$ (240) & days $0$--$44$ (45) & days $45$--$120$ (76)\\
\bottomrule
\end{tabular}
\end{table}

\subsection{Training recipe}
\label{sec:recipe}

Table~\ref{tab:recipe} lists every setting shared by the Lotka-Volterra
runs; each experiment in Section~\ref{sec:results} changes exactly one of
them.

\begin{table*}[t]
\centering
\footnotesize
\renewcommand{\arraystretch}{1.1}
\caption{Shared training recipe for all Lotka-Volterra runs.}
\label{tab:recipe}
\begin{tabularx}{\textwidth}{@{}L{3.6cm} L{5.2cm} Y@{}}
\toprule
\hdrrow \thh{Setting} & \thh{Value} & \thh{Note}\\
\midrule
Network & KAN $[2,10,2]$, $G=5$, 240 parameters & knots uniform on $[-1,1]$, $h=0.5$\\
Edge function & Equation~\ref{eq:kan-edge}, $\tanh$ normaliser, SiLU residual & Glorot-uniform initialisation of $\mathbf{C}$ and $\mathbf{W}$\\
Integrator & Tsit5, fixed step & 2 sub-steps per sample: $h=\Delta t/2=0.05$\\
Sampling interval & $\Delta t=0.1$ & 36 training points, 141 on $[0,14]$\\
Loss & mean squared trajectory error on $[0,3.5]$ & no regularisation ($\lambda=0$ in Equation~\ref{eq:l1-sparsity})\\
Optimiser & Adam \cite{kingma2015adam}, learning rate $2\times10^{-3}$ & constant learning rate, no schedule\\
Gradient clipping & global norm $\le1.0$ & non-finite gradients skip the step (Section~\ref{sec:crossdomain})\\
Budget & $10{,}000$ epochs, full batch & $50{,}000$ in the budget validation (Section~\ref{sec:phase4})\\
Model selection & checkpoint with minimum training MSE & the test window is never used for selection\\
Precision, seed & float32, seed 42 & one seed per configuration ($N=1$)\\
\bottomrule
\end{tabularx}
\end{table*}

\subsection{Evaluation protocol}
\label{sec:protocol}

Model selection uses the minimum \emph{training}-window loss only, and every
accuracy number is computed strictly outside the training horizon
(Figure~\ref{fig:data-split}), so a number cannot improve just by widening
what counts as fit.

\begin{figure*}[t]
\centering
\includegraphics[width=0.85\textwidth]{figures/tikz/data_split.pdf}
\caption{Time axis of the Lotka-Volterra protocol. Tick marks above the bar
are multiples of the true period. The far-horizon window $(14,28]$ is used
only for the stress tests of Sections~\ref{sec:crossdomain}
and~\ref{sec:phase4}; no model is retrained or reselected for it.}
\label{fig:data-split}
\end{figure*}

Writing $\hat{\mathbf{u}}_k$ for the model's prediction and $\mathbf{u}_k$
for the truth at the $N_e=105$ test samples, the reported metrics are
\begin{gather}
\mathrm{MSE}=\frac{1}{2N_e}\sum_{k}\|\hat{\mathbf{u}}_k-\mathbf{u}_k\|_2^2,\qquad
R^2 = 1-\frac{\sum_k\|\hat{\mathbf{u}}_k-\mathbf{u}_k\|_2^2}{\sum_k\|\mathbf{u}_k-\bar{\mathbf{u}}\|_2^2},
\nonumber\\
\text{rel.}\,L_2=\frac{\big(\sum_k\|\hat{\mathbf{u}}_k-\mathbf{u}_k\|_2^2\big)^{1/2}}{\big(\sum_k\|\mathbf{u}_k\|_2^2\big)^{1/2}},
\label{eq:metrics}
\end{gather}
where $\bar{\mathbf{u}}$ is the componentwise mean of the true test
trajectory. $R^2=1$ is a perfect match, $R^2=0$ is no better than
predicting that mean, and a negative value is worse than that trivial
baseline. Two further quantities describe the trained field rather than its
trajectory. The \textbf{Lipschitz estimate}
\begin{equation}
\hat L=\max_{k\le 36}\frac{\|\mathbf{f}_\theta(\mathbf{u}_k+\varepsilon\mathbf{v}_k)-\mathbf{f}_\theta(\mathbf{u}_k)\|_2}{\varepsilon},
\label{eq:lipschitz}
\end{equation}
with $\varepsilon=10^{-3}$ and random unit vectors $\mathbf{v}_k$, is a
finite-difference sample of the local slope of $\mathbf{f}_\theta$ at the
training states. It is a \emph{lower} bound on the field's true Lipschitz
constant and measures how sharply the learned field bends. The
\textbf{cost} of a solver is its number of function evaluations (NFE): an
$s$-stage method on the training window costs $\text{NFE}=35\times2\times
s=70s$ calls of $\mathbf{f}_\theta$ per epoch. Because reverse-mode
differentiation unrolls every call, training wall-clock scales with it too.

\subsection{Computational budget: why 10{,}000 epochs}
\label{sec:budget}

The base paper trains for $10^5$ epochs. Table~\ref{tab:cost} shows what
that would cost us. The paper's integration loop is compiled by Julia's
SciML stack; ours is interpreted PyTorch on CPU, where every Runge-Kutta
stage is a separate Python-level call with its own autograd bookkeeping. On
our development machine one Tsit5/RBF epoch takes $0.17$~s and one
Tsit5/B-spline epoch $0.67$~s; the cloud CPUs we used for the long runs of
Section~\ref{sec:phase4} were $2.2$--$2.5\times$ slower for every KAN
configuration.

\begin{table*}[t]
\centering
\footnotesize
\renewcommand{\arraystretch}{1.1}
\caption{Measured per-epoch cost on our local development CPU and on cloud
(Kaggle) CPU, each averaged over the full run, and the projected cost of the
paper's $10^5$-epoch budget.}
\label{tab:cost}
\begin{tabular}{@{}L{4.2cm} C{1.6cm} C{1.6cm} C{1.2cm} C{2.5cm} C{2.5cm}@{}}
\toprule
\hdrrow \thh{Configuration} & \thh{Local s/epoch} & \thh{Cloud s/epoch} & \thh{Ratio} & \thh{$10^5$ epochs, local} & \thh{$10^5$ epochs, cloud}\\
\midrule
Euler / RBF & $0.0254$ & $0.0627$ & $2.47\times$ & $0.7$ h & $1.7$ h\\
Tsit5 / RBF (reference) & $0.1705$ & $0.4033$ & $2.37\times$ & $4.7$ h & $11.2$ h\\
Tsit5 / B-spline & $0.6680$ & $1.5438$ & $2.31\times$ & $18.6$ h & $42.9$ h\\
MLP-ODE, SiLU $[2,14,8,8,2]$ & $0.2856$ & $0.2722$ & $0.95\times$ & $7.9$ h & $7.6$ h\\
MLP-ODE, tanh $[2,50,2]$ & $0.0736$ & $0.1595$ & $2.17\times$ & $2.0$ h & $4.4$ h\\
\midrule
All 24 Lotka-Volterra configurations & \multicolumn{2}{c}{$14.2$ CPU-h at $10^4$ epochs} & & $142$ h & $\approx336$ h\\
\bottomrule
\end{tabular}
\end{table*}

A $10^5$-epoch version of the 24-configuration Lotka-Volterra sweep would
need roughly $142$ CPU-hours locally, or about two weeks of cloud CPU time,
before the pendulum, SIR and follow-on studies. We therefore ran the sweep
at $10^4$ epochs, a budget at which every noise-free solver and basis
configuration except Euler and four bases (Lagrange, IQF, RSWAF, Newton)
reaches $10^{-4}$ training MSE (Table~\ref{tab:milestones}). That choice is
only safe if the conclusions drawn at $10^4$ epochs survive a longer budget,
which is exactly what Section~\ref{sec:phase4} tests.
